\input{../tex-inputs/latex-leseansicht-vorspann}

\section[Arthur Schnitzler an Gustav Schwarzkopf, {[}6. oder 7{]}. 3. 1905]{L04393 Arthur Schnitzler an Gustav Schwarzkopf, {[}6. oder 7{]}. 3. 1905}
\nopagebreak\mylabel{L04393v}
\rehead{ }\normalsize\beginnumbering\briefempfaengerindex{Schwarzkopf, Gustav@\textsc{Schwarzkopf, Gustav}!zzzSchnitzler, Arthur@\emph{von Arthur Schnitzler}!1905-03-072@{{[}6. oder 7{]}. 3. 1905}|(be}
\toendnotes[C]{\smallbreak\pagebreak[2]}
\correspDesc{Versand  durch Arthur Schnitzler im Zeitraum [6. oder 7]. 3. 1905 in Ajaccio
\newline{}Übermittlung  am 7. 3. 1905 in Ajaccio
\newline{}Erhalt  durch Gustav Schwarzkopf im Zeitraum [9. 3. 1905 – 14. 3. 1905?] in Tiefer Graben 23}\toendnotes[C]{\smallbreak}
\Standort{DLA, A:Schnitzler, HS.1985.1.1897.}
\physDesc{Bildpostkarte, 129 Zeichen
\newline{}Handschrift: Bleistift, lateinische Kurrent
\newline{}Versand: Stempel: »\nobreak{}\oindex{Ajaccio@\textbf{Ajaccio}|pwk}Aj{[}accio{]} Corse, 7–3 05, \textcolor{gray}{H} 10\nobreak{}«.  }\pstart{}{\pb}Autriche\oindex{Österreich@\textbf{Österreich}|pw}\pend{}\pstart{}\textcolor{gray}{\textbf{M}} Gustav Schwarzkopf\pend{}\pstart{}Wien I\oindex{I., Innere Stadt@\textbf{I., Innere Stadt}|pw}.\pend{}\pstart{}Tiefer Graben 23\oindex{Wien@\textbf{Wien}!I., Innere Stadt@\textbf{I., Innere Stadt}!Tiefer Graben 23@\textbf{Tiefer Graben 23}|pw}.\pend{}{\bigskip}
\pstart
           \centering{}{\pb}\textcolor{gray}{\textbf{Ajaccio\oindex{Ajaccio@\textbf{Ajaccio}|pw}. – Arrivée du Courrier.}}\pend
           \vspace{1em}
\pstart
           \noindent{}{\pb}Herzliche Grüsse! Es ist wunderschön, echter Frühling! Grüssen Sie Dr. Max\pwindex{Schwarzkopf, Max 12.\,6.\,1857 Wien – 14.\,4.\,1928 ebd.@\textsc{Schwarzkopf, Max} (12.\,6.\,1857 Wien – 14.\,4.\,1928 ebd.)|pw}.\pend
           
\pstart
           Ihr{\\[\baselineskip]}\spacefill\mbox{A.}\pend
           \leftskip=0em{}\selectlanguage{ngerman}\endnumbering\briefempfaengerindex{Schwarzkopf, Gustav@\textsc{Schwarzkopf, Gustav}!zzzSchnitzler, Arthur@\emph{von Arthur Schnitzler}!1905-03-062@{{[}6. oder 7{]}. 3. 1905}|)be}\mylabel{L04393h}
\begin{anhang}
\end{anhang}\newcommand{\dateiname}{L04393}\newcommand{\titel}{Arthur Schnitzler an Gustav Schwarzkopf, [6. oder 7]. 3. 1905}\newcommand{\editorInnen}{Herausgegeben von Jahnke, SelmaMüller, Martin Anton}\input{../tex-inputs/latex-leseansicht-abspann}
